% page 528 of the facsimile (image p-527 of the scan)
% block 154.9 x 237.7 mm ; left margin 8.0 mm
\pgc{-12.700mm}{0.000mm}{
\l{\cel{3}520.}
\l{}
\l{\sou{sero}. (\sou{patol}.) I. Liquido aquatra qua su separas de la}
\l{\cel{3}parto koagulinta di sango (o di lakto). - II. (\sou{fiziol}.)}
\l{\cel{3}Humoro sekrecita (da plevro, peritoneo, e c.) - III.}
\l{\cel{3}(\sou{patol.}) Humoro ne-normala, en ula morbi. - DEFIRS.}
\l{}
\l{\sou{seraf}o. (\sou{teol.}) Anjelo di la hierarkio unesma ; spirito}
\l{\cel{3}cielala, che la Judi e che la kristani. - DEFIS.}
\l{}
\l{\sou{serayo}. Palaco di la sultani, viziri, siniori Turka (en}
\l{\cel{3}Turkia olima). - DEFIRS.}
\l{}
\l{\sou{serbatano}. I. Tubo longa, kava, uzata por lansar, per homo-}
\l{\cel{3}suflo, pizi, mikra buli ek tero, e c. - II. Tubo de fero,}
\l{\cel{3}uzata por ensuflar vitro. - FIS.}
\l{}
\l{\sou{serchar}.\cel{2}(\sou{trans.)}\cel{3}Probar deskovror, prokuror a su, ren-}
\l{\cel{3}kontror (ulu, ulo). - EFI.}
\l{}
\l{serena. (aludante cielo, atmosfero, e c., o psiko).Pura e}
\l{\cel{3}kalma, quieta. - EFIS.}
\l{}
\l{\sou{serenado}.Koncerto, donata vespere sub la fenestri di ulu}
\l{\cel{3}(precipue en Hispania), olim. - DEFIRS.}
\l{}
\l{\sou{serfo.(feudismo-rejimo}).\cel{2}Homo pri la persono o la havaji}
\l{\cel{3}di qua la sinioro havis yuri qui variis segun la kus-}
\l{\cel{3}tumi. - EFIS.}
\l{}
\l{\sou{serio}. I. Termini qui intersequas segun lego (+\sou{leyo}) deter-}
\l{\cel{3}minita. - II. Intersequentaro kontinua. - DEFIRS.}
\l{}
\l{\sou{serino}. (\sou{zool.}) Pasero di qua la plumaro esas citron-flava}
\l{\cel{3}che plura speci, e di qua la maskulo povas lernar siflo}
\l{\cel{3}e kantar ula arii. - L.\cel{1}\sou{fringilla serinus}. - EF.}
\l{}
\l{\sou{seringo}. (\sou{bot.}) Arboreto tre rametoza, kun folii inter-opo-}
\l{\cel{3}zanta, kun flori blanka qui dispozesas bukete, e tre}
\l{\cel{3}odoras. - L.\cel{1}\sou{philadelphus.}\cel{2}- EFIS.}
\l{}
\l{\sou{serioza}. Qua egardas, che la kozi, lia gravajo, lia importoz-}
\l{\cel{3}ajo, prefere. - EFIRS.}
\l{}
\l{\sou{serjo}. Stofo de silko, de filo lina, de lano, di qua la}
\l{\cel{3}wefto esas min glata e min densa kam la fili longesala}
\l{\cel{3}(warpa), qui esas apene videbla averse. - DEF.}
\l{}
\l{\sou{serjento}. I. Oficiro qua resortisas a la autoritatozi sta-}
\l{\cel{3}tala, municipala od urbala, e qua havis plura roli ple-}
\l{\cel{3}enda. - II. (\sou{armeo)}. Suboficiro di infantrio o di kavalrio}
\l{\cel{3}DEFIRS.}
\l{}
\l{\sou{serono}. Pakego, envelopita per bovo-felo (e qua kontenas}
\l{\cel{3}vari exotika). - DEFIS.}
}
